\section{Protocols and extended experimental evidence}
\label{app:extended}
\subsection{Experimental units and statistics}
\label{app:protocols}
Table~\ref{tab:studies} separates the study units. A foundation is a sensory
checkpoint; an action seed initializes or trains its action path. Multiple
action seeds under one foundation are not treated as independent foundations.
\begin{table}[tbp]\centering\small
\caption{Experimental units and the question answered by each study.}
\label{tab:studies}
\begin{tabular}{@{}p{.24\linewidth}p{.31\linewidth}p{.39\linewidth}@{}}\toprule
Study & Units & Question\\\midrule
Main stream & 3 foundations, 3 action seeds each & Prediction, tracking, storage, return correction, latency\\
Routed actuator RLS & Same $3\times3$ layout, separate replay & Actuator-space identification under the six-stage protocol\\
Component roles & 5 model-seed clusters & Matrix contribution versus score-only and no-inverse variants\\
Behavioral equivalence & 4 controllers, 3 tasks, 3 environment seeds & Bounded versus legacy history; first changed command\\
Closed loop & 5 models, 3 tasks, 5 replicas per arm & Reward and success under attenuation\\
Joint learning & 5 action seeds, one sensory foundation & Shared-core predictive and semantic learning\\
Training-stage pilot & One foundation; matched cached stream & Same adaptation interface across offline stages\\
Controlled probes & 3 seeds per separate model study & Temporal, future-statistic, and factor-association tests\\
Output calibration & One foundation; 4 paired units per arm & Hidden versus output-space correction\\\bottomrule
\end{tabular}\end{table}

For the main stream, each action-seed run first averages the measured shifted
transition errors, then computes its gain relative to the paired frozen arm.
There are 126 measured shifted transitions per run, after omitting each
episode's four-transition prefix from the score. All methods replay the same
observations; tracking is evaluated through the known actuator map.
Action seeds are averaged within each foundation, and the reported mean
averages foundations. The archived 95\% intervals use 20,000 bootstrap draws
of foundation means. Update latency uses 2,988 synchronized samples per
adaptive arm after discarding 16 warm-up updates from each run.

Source stages reuse router-calibration episodes and four-transition prefixes;
the provenance record identifies 48 reused source-prefix rows per evaluated
replica. Shifted stages use a separate execution archive. This defines the
calibrated-source reentry measurement in the main table. The sequence can
retain matrices and scores across stages while re-establishing the route at
the start of each episode.

Closed-loop statistics preserve method pairing while resampling model, task,
and replica axes independently, with 20,000 draws. Success is the episode-level
indicator of task completion; cumulative reward sums environment rewards.
The reference code's older control runner has its own flat-episode bootstrap;
the released recorded results contain the multi-axis intervals reported here.

\subsection{Independent foundation results}
\begin{table}[tbp]
\centering\small
\caption{Prediction gain by independent foundation (\%). Each row averages three action seeds.}
\label{tab:foundations}
\begin{tabular}{@{}lrrrr@{}}
\toprule
Foundation & External state & Adapter & Prediction-path & Replay \\
\midrule
33401 & 21.80 & 14.60 & 25.30 & 25.08 \\
33402 & 20.34 & 17.42 & 17.47 & 16.76 \\
33403 & 15.71 & 14.82 & 22.79 & 19.29 \\
\bottomrule
\end{tabular}
\end{table}

The three foundation rows show how prediction gain varies across independently
trained foundations. The frozen reference has zero gain by definition. The main result is a
mean trade-off over these independently trained foundations; all rows and
comparison arms are retained in the evidence files.

The same-protocol routed actuator RLS uses 4,112 bytes, including its routing
prefix. It obtains 0.58\% prediction gain and $-1.05$\% tracking gain; the
replicated external-state arm obtains 19.72\% and 3.87\%, respectively.
Both return zero correction on calibrated source revisits. This RLS estimates
actuator coupling and is distinct from the output-space calibrators below.

\subsection{Controlled temporal and cross-modal probes}
\label{app:probes}
These historical studies use separate event encoders, with predictive and
contrastive objectives, and are not additional main-stream foundations.
Identity metadata builds paired views; the future-prediction targets themselves
are computed from observations.

In the temporal probe, the full sequence is compared with two controls.
Order shuffling permutes non-time event features across fixed timestamps.
The rate-grid control sorts event features and assigns evenly spaced
timestamps. The outcome is future-slice negative log-likelihood (NLL), where
lower values indicate better prediction. These controls test sensitivity to
the temporal structure of this task.

The future-statistic probe predicts normalized histograms relating event
time and feature coordinates in later image/audio views. A learned encoder
is compared with a random encoder using the same prediction task. The
composition probe receives factor segments and compares true image/audio
factor associations with shuffled associations. Its score is top-1 factor
association accuracy. Table~\ref{tab:probes} reports the archived three-seed
means; the separate objectives have different metric scales.
\begin{table}[tbp]
\centering\small
\caption{Controlled event-model probes, three-seed means. Lower NLL/loss and higher accuracy are better. Temporal, future-statistic and factor-association studies correspond to archived M60, M69 and M71 records.}
\label{tab:probes}
\begin{tabular}{@{}lllr@{}}
\toprule
Study & Condition & Metric & Mean \\
\midrule
Temporal & Ordered sequence & NLL & 4.9118 \\
Temporal & Rate-grid control & NLL & 6.9293 \\
Temporal & Shuffled order & NLL & 11.0357 \\
Future statistics & Learned encoder & Loss & 0.1103 \\
Future statistics & Random encoder & Loss & 0.1397 \\
Factor association & True pairing & Accuracy (\%) & 46.6667 \\
Factor association & Shuffled pairing & Accuracy (\%) & 6.6667 \\
\bottomrule
\end{tabular}
\end{table}


\subsection{Three-stage joint-learning results}
\label{app:joint}
Table~\ref{tab:joint_stages} reports all three stages of the five-action-seed
study. Its sensory metric averages four groups: audio (the mean of SHD and
SSC), image, text, and video. The stored normalized values are used as
recorded. The reference evaluator separately uses validation-set persistence
denominators and an unweighted five-stream mean; that evaluator's aggregate
should not be substituted for this archived four-group statistic.
Relative to action-only updates, joint training increases action MSE by
3.30\% on average across paired seeds (95\% CI: 1.65\%--4.94\%), while
improving class retrieval by 3.31 percentage points (95\% CI: 0.94--5.38).
Joint updates improve retrieval, while action-only updates favor action
prediction.
\begin{table}[tbp]
\centering\small
\caption{All stages of the five-action-seed joint-learning study. Sensory MSE is the archived four-group normalized mean; accuracy and retrieval are percentages.}
\label{tab:joint_stages}
\begin{tabular}{@{}lrrrr@{}}
\toprule
Stage & Action MSE & Sensory MSE & Semantic (\%) & Retrieval (\%) \\
\midrule
Sequential & 0.1463 & 0.6776 & 19.1488 & 10.0000 \\
Action-only shared update & 0.1174 & 0.6786 & 19.2655 & 8.5000 \\
Joint shared update & 0.1212 & 0.6725 & 20.9440 & 11.8125 \\
\bottomrule
\end{tabular}
\end{table}


\subsection{Training-stage comparison}
\label{app:bridge}
The pilot begins with one common sensory-pretrained checkpoint. From the
sequential action-trained model, three copies receive 200 additional updates:
joint training, action loss alone, or joint training without the sensory
next-slice term. Thus the no-next variant removes the term in the final stage,
not from the full pretraining history. Offline fitting supplies each model's
basis and router using the same procedure and hyperparameters.

All models replay the same cached six-stage trajectory and compare frozen,
score-only, and full external state on 126 shifted transitions after routing.
Table~\ref{tab:bridge} shows that the interface can be used across these
training stages. The ordering of adapted errors differs from the ordering of offline action
errors. Lower offline action error alone does not determine which checkpoint
has the lowest error after adaptation in this pilot.
\begin{table}[tbp]
\centering\small
\caption{One-foundation training-stage pilot. MSE is evaluated on 126 shifted transitions after the observation prefixes. All variants inherit sensory predictive pretraining.}
\label{tab:bridge}
\begin{tabular}{@{}lrrrr@{}}
\toprule
Training stage & Frozen MSE & Score MSE & Full MSE & Gain (\%) \\
\midrule
Sequential & 0.3927 & 0.3661 & 0.2878 & 26.72 \\
Action-only & 0.4408 & 0.3797 & 0.2918 & 33.81 \\
Joint & 0.4097 & 0.3992 & 0.3116 & 23.96 \\
Joint, no-next & 0.4270 & 0.4129 & 0.3186 & 25.39 \\
\bottomrule
\end{tabular}
\end{table}


\subsection{Direct output-space calibration}
\label{app:additional}
The additional development comparison applies corrections at the output
rather than the hidden feature. Table~\ref{tab:additional} preserves all
methods: hidden residual state, direct-output NLMS, direct-output RLS, and a
low-rank readout update. It uses one sensory foundation and four paired units
per method, rather than the three-foundation main comparison.
\begin{table}[tbp]
\centering\small
\caption{Separate output-calibration development protocol: one sensory foundation, four paired units per method. These results are not pooled with the main-stream comparison.}
\label{tab:additional}
\begin{tabular}{@{}lrrrr@{}}
\toprule
Correction & MSE & Gain (\%) & Bytes & Update (ms) \\
\midrule
Hidden residual state & 0.2640 & 23.34 & 28,848 & 11.70 \\
Output NLMS & 0.1652 & 51.70 & 13,956 & 11.03 \\
Output RLS & 0.1839 & 45.84 & 21,156 & 11.10 \\
Low-rank readout & 0.2961 & 14.14 & 30,532 & 11.97 \\
\bottomrule
\end{tabular}
\end{table}


In this setting, direct-output NLMS and RLS achieve larger prediction gains
with less mutable state than the hidden-state correction. This comparison favors output-space correction in this development setting.
We report it separately from the main-stream frontier because the protocols
differ. The released records include every method arm with its protocol.

\subsection{Closed-loop detail}
\label{app:control}
Each model/task controller is initialized once, then processes five fresh
trajectory replicas in order. External state resets its route at each episode
but keeps its matrices and candidate scores. Scalar RLS and gain-memory replay
also retain their accumulated state across these replicas. All episodes use
attenuation 0.4 and last 13 decisions, or 104 simulator steps.

The live-control tuning arm uses AdamW at learning rate $10^{-5}$ on parameters
whose names do not contain \texttt{semantic}. Gain-memory replay retrieves
local gain estimates from observation/task features using nearest neighbors,
with capacity 141. It does not train the prediction path with replay as the
main-stream Adam arm does. The scalar RLS controller accumulates a numerator
and denominator using finite-difference model sensitivity. The oracle divides
by the true attenuation and clips to the valid command range; unlike the
external-state controller, it has no 0.05 correction cap. Consequently, the
main-stream storage and latency entries are not assigned to these live-control
configurations.

All six arms appear in Table~\ref{tab:control}, including the scalar RLS and
oracle references. The external-state reward gains by task are 2.31 for
pick-place, 11.61 for push, and 9.79 for reach. The five model-seed gains range
from 7.19 to 8.66. Prediction-model tuning and external state both achieve
66.67\% aggregate success. The scalar RLS obtains the largest reward among
the non-oracle methods in this single-fault control protocol. External state improves reward over frozen control, while scalar RLS is the
stronger controller for this particular fault.
